\subsection{对偶容量形式的修正与强对偶}\label{sec:r17:dual}

把 $\kappa_3$ 写成「方向组 $U$ 固定、只对质量 $g$ 优化」的形式，会得到一个干净的对偶。

\begin{lemma}[对偶容量]\label{lem:phid}
对固定的 $U=(u_i)\subset S^2$ 与 $c\in\R^N$，令
$\Phi_U(c):=\inf\bigl\{\sum_ic_i^2/g_i:\ g>0,\ \sum_ig_iu_iu_i^\top=I_3\bigr\}$。则
\begin{equation}\label{eq:phid}
\Phi_U(c)=\max_{M\,:\,u_i^\top Mu_i\ge0\ \forall i}\ \frac{\bigl(\sum_i|c_i|\sqrt{u_i^\top Mu_i}\bigr)^2}{\tr M},
\end{equation}
右端是\textbf{尺度不变}的（$M\to tM$ 不变），因而是紧集 $\{M\succeq0,\tr M=1\}$ 上的凹函数最大化。
\end{lemma}

\begin{proof}
对任意满足 $u_i^\top Mu_i\ge0$ 的 $M$，令 $\psi_i=\sqrt{u_i^\top Mu_i}$。
AM--GM 给 $c_i^2/g_i+g_i\psi_i^2\ge2|c_i|\psi_i$；对 $i$ 求和并用帧条件
$\sum_ig_i\psi_i^2=\tr\bigl(M\sum_ig_iu_iu_i^\top\bigr)=\tr M$，得
$\sum c_i^2/g_i\ge2\sum|c_i|\psi_i-\tr M$。
对固定的方向 $M_0$ 再优化标度 $M=tM_0$（$2\sqrt tA-tB$ 的最大值是 $A^2/B$）即得 \eqref{eq:phid} 的「$\ge$」；
「$\le$」由 KKT：最优 $g_i=|c_i|/\psi_i$ 使 AM--GM 逐点取等。
\end{proof}

\begin{corollary}[$\kappa_3$ 的容量表述]\label{cor:phikappa}
$\kappa_3=\sup_{U,\,c\ne0}\dfrac{c^\top K_+(U)c}{\Phi_U(c)}$，其中 $K_+(U)_{ij}=(\langle u_i,u_j\rangle)_+$。
故 $\kappa_3\le\rho_0$ 等价于：对一切 $U$ 与 $c$，$c^\top K_+(U)c\le\rho_0\,\Phi_U(c)$。
\end{corollary}

\begin{proof}
$\lambda_{\max}(D_g^{1/2}K_+D_g^{1/2})=\max\{c^\top K_+c:\sum_ic_i^2/g_i=1\}$，再对 $g\in\mathcal G_U$ 取上确界。
\end{proof}

冠军处的数值：最优 $M=\Lam$（正是定理~\ref{thm:stiefelkkt} 的迹一乘子），$\tr M=1$，
$\psi_i=\sqrt{u_i^\top\Lam u_i}=h_i$，$\Phi=1$，$c^\top K_+c=\rho_0$，
\textbf{零对偶间隙}（primal $=0.999999978$，dual $=1.000000000$，差 $2.2\times10^{-8}$；
由 $g_i=c_i/\sqrt{t\,u_i^\top Mu_i}$ 复原的质量与 primal 最优质量之差 $2.25\times10^{-8}$）。
在 \texttt{random\_frame\_measure} 给出的可行构型（$N=5,\dots,8$）上，随机方向 $c$ 的强对偶相对差 $\le1.3\times10^{-4}$。

\begin{remark}[一处必须记录的可行性陷阱]\label{rmk:phitrap}
若 $U$ 的原子使 $I_3$ \textbf{不}在 $\{\sum_ig_iu_iu_i^\top:g\ge0\}$ 的内部（例如 $6$ 个随机方向），
则 primal 不可行、$\Phi_U(c)=+\infty$，对偶 \eqref{eq:phid} 也发散到 $+\infty$
（脚本上一度出现 $1.6\times10^{17}$）。这\textbf{不是}强对偶的反例，而是可行域为空。
上一版脚本还把 primal 目标误写成 $\sum c_i^2g_i$（应为 $\sum c_i^2/g_i$），
造成 $\Phi_{\rm primal}=0.2503$ 的假违背；两处均已修正，修正后 \S(F) 全绿。
\end{remark}

\subsection{钝图的团数与 Motzkin--Straus}\label{sec:r17:ms}

\begin{lemma}[钝团数]\label{lem:obtuseclique}
钝图 $G_-=(\{1,\dots,N\},\{\{i,j\}:\langle u_i,u_j\rangle<0\})$ 满足 $\omega(G_-)\le4$，且 $4$ 可达。
\end{lemma}

\begin{proof}
$\R^3$ 中两两成钝角的单位向量至多 $4$ 个（若 $u_1,\dots,u_5$ 两两钝，
则 Gram 矩阵 $G$ 满足 $G\mathbf 1$ 的各分量 $=1+\sum_{j\ne i}\langle u_i,u_j\rangle$，
而 $\rank G\le3$ 使 $G$ 有一个零空间向量；标准论证给 $k\le d+1$）。
可达性：数值扫描中确实出现 $\omega(G_-)=4$ 的构型（$360$ 个随机构型的最大值恰为 $4$）。
\end{proof}

\begin{corollary}[Motzkin--Straus 型界]\label{cor:ms}
记 $\varpi:=\sum_{\{i,j\}\text{ 钝}}x_i^2x_j^2$（无序对求和的两倍，即 $\sum_{i\sim j}x_i^2x_j^2$）。则
$\varpi\le\frac12\bigl(1-\frac1{\omega(G_-)}\bigr)\le\frac38$，且 $\beta\le\sqrt{e\,\varpi}\le\sqrt{3e/8}$。
\end{corollary}

\begin{proof}
以 $y_i=x_i^2$（$\sum y_i=1$）代入 Motzkin--Straus 定理；
$\beta^2\le e\varpi$ 是把 $\beta=\sum_{ij}(Q_-)_{ij}x_ix_j$ 与 $\|Q_-\|_F^2=e$、$\varpi$ 作 Cauchy--Schwarz。
\end{proof}

数值：$360$ 个构型上 $\max\varpi=0.359148483447<3/8=0.375$，$\beta^2-e\varpi\le0$ 无一违背。

\begin{remark}[这条路线的定量天花板（诚实记录）]\label{rmk:mslimit}
把 $\beta\le\sqrt{3e/8}$、$\rho^2\le\|Q_+\|_F^2=3-e$（引理~\ref{lem:frob}）与 $\tau\le\beta$ 联立：
$\tau\le\sqrt{3(3-(1+\tau)^2)/8}$，解得 $\tau\le0.5146$、$\rho\le1.5146$，
\textbf{比 Mohar 的 $1.36603$ 还差}。原因见注~\ref{rmk:taugap}：
Motzkin--Straus 只控制 $\beta$，而\textbf{抵消项 $-(1-\alpha)$ 完全没有被用上}。
本小节因此记为「已证但不足」，并作为「必须让 $(I-Q)x$ 参与」这一判断的定量佐证。
\end{remark}

\subsection{本轮状态、诚实标签与方法优先队列}\label{sec:r1517:status}

\textbf{已证（符号推导 $+$ 数值证书）}：
命题~\ref{prop:gradformula}（正部的环境梯度，勘误后重新确立）；
定理~\ref{thm:stiefelkkt}（Stiefel KKT 与迹一刚性）；
推论~\ref{cor:ellipsoid}（椭球面恒等式）、\ref{cor:hubfree}（枢纽无关式 $\rho^2=m^\top\Lam^{-1}m$ 与精确 CS 缺陷）；
引理~\ref{lem:tauform}（$\tau$-形式）与推论~\ref{cor:oneneg}、引理~\ref{lem:normqmx}；
引理~\ref{lem:frob}（能量守恒 $\|Q_+\|_F^2=3-e$，精确）、推论~\ref{cor:floor2}（一行式谱地板）；
引理~\ref{lem:framesplit}、\ref{lem:cs1}、\ref{lem:cs2}、定理~\ref{thm:betaAA}、命题~\ref{prop:betaAAeq}；
引理~\ref{lem:firstorder}、定理~\ref{thm:secondvar2}、\textbf{定理~\ref{thm:dichotomy}（全局二分）}、
推论~\ref{cor:branch2rigid}（分支 II 的球面二次曲线刚性）；
引理~\ref{lem:phid} 与推论~\ref{cor:phikappa}（对偶容量、强对偶）；
引理~\ref{lem:obtuseclique} 与推论~\ref{cor:ms}（$\omega(G_-)\le4$、$\varpi\le3/8$、$\beta\le\sqrt{3e/8}$）。

\textbf{条件性（已证但需附加假设）}：
命题~\ref{prop:majbound}（$Q_+\succeq0$ 且 $n_+\le4$ 时 $\rho\le1.290244<\Phi$）；
推论~\ref{cor:ellipsoid} 的 $\Lam\succ0$ 要求（冠军处成立：谱 $\{0.5946,0.2027^{[2]}\}>0$）。

\textbf{否证（同样是成果）}：
命题~\ref{prop:fourthfalse}——全局第四不等式 $\alpha+\sqrt{A_1A_2}\le\rho_0$ 为\textbf{假}，
$N=7$ 反例超出 $6.27\times10^{-3}$；
注~\ref{rmk:mslimit}——Motzkin--Straus 路线的定量天花板 $1.5146$，不足；
注~\ref{rmk:majcond}——重排路线在 $n_+=6$（冠军的实际层）只给 $1.474$，不足。

\textbf{数值（不作成果）}：
$N=5,\dots,12$ 的 Stiefel 梯度上升极大元及其质量分组；
分支 II 在 $N=12$ 极大元上的精确实现（$R|_{\mathcal K}=0$ 到 $3\times10^{-15}$）；
$360$ 个随机构型上的能量不等式扫描。

\textbf{方法优先队列（第七版）}
\begin{enumerate}\itemsep3pt
\item \textbf{P1（队首，全局）——排除分支 II}：证明推论~\ref{cor:branch2rigid} 的插值刚性
\eqref{eq:conic} 在原子两两不重合时不可能成立（$N\ge7$）。
这是一道\textbf{有限维球面插值}题：一个二次多项式要在 $N$ 个点上等于
$\sqrt{u^\top\Lam u}\cdot(\langle u_i,u\rangle)_+$，后者的折点结构是主要武器。
\item \textbf{P2（全局）——压低 $n_+$ 到 $4$}：推论~\ref{cor:floor2} 说 $n_+\ge4$，
定理~\ref{thm:dichotomy}(i) 说 $n_+\le7$，命题~\ref{prop:majbound} 说 $n_+=4$ 即可完成「优于 \STW」。
缺口是 $n_+\in\{5,6,7\}$ 三层（冠军在 $n_+=6$）。
\item \textbf{P3（局部）——分支 I}：在 $\nu_-(Q_+)\ge1$ 的构型上直接证 $\rho\le\rho_0$。
可用 $\nu_-(Q_+)\le\nu_-(Q_-)$ 与 $\omega(G_-)\le4$（引理~\ref{lem:obtuseclique}）。
\item \textbf{P4（$N=6$ 收尾）}：第~\ref{sec:n6global} 节余下的 $20$ 个符号类。
\item \textbf{P5（副方向，下界）}：$\Lambda_3(G)\ge\rho_0n-3$ 的可达到性——冠军是 hub$+F_5$，
对应的图族已由定理~\ref{thm:Fn} 给出闭式谱，下界一侧只差「无穷族的渐近」这一步。
\end{enumerate}

\textbf{本轮把「唯一缺口」重写为}：定理~\ref{thm:dichotomy} 的分支 I 与分支 II 二者之一必须被排除。
分支 II 的形式已完全代数化（\eqref{eq:conic}），分支 I 的形式是一个负惯性指数条件。
两者都不再需要「整条矩核直线上 $Q_+\succeq0」这种难以处理的断言。

